% ============================================================
%   SESSION 2025
% ============================================================

\sujetbacheader{2025}


% ------------------------------------------------------------
%   EXERCICE 1
% ------------------------------------------------------------

\exercice{3 points}

\noindent\textbf{Partie A}

Soient $A$ et $B$ deux matrices carrées d'ordre $3$ définies par :
$$
A = \begin{pmatrix} 0 & 1 & -1 \\ 1 & 0 & 3 \\ 1 & 1 & 1 \end{pmatrix}
\quad \text{et} \quad
B = \begin{pmatrix} -3 & -2 & 3 \\ 2 & 1 & -1 \\ 1 & 1 & -1 \end{pmatrix}.
$$

\begin{enumerate}
	\item
	\begin{enumerate}
		\item Calculer $A \times B$ et $B \times A$. \hfill (0,25 pt $\times$ 2)
		\item Que peut-on en conclure pour les deux matrices $A$ et $B$ ? \hfill (0,25 pt)
	\end{enumerate}
	\item On considère le système $(S)$ :
	$$
	\begin{cases}
		y - z = -3 \\
		x + 3z = 17 \\
		x + y + z = 11
	\end{cases}
	$$
	\begin{enumerate}
		\item Écrire ce système sous forme matricielle : $AX = Y$. \hfill (0,25 pt)
		\item Déterminer la matrice $X$. \hfill (0,5 pt)
	\end{enumerate}
\end{enumerate}

\bigskip

\noindent\textbf{Partie B}

\begin{enumerate}
	\item Déterminer le reste de la division de $7 \times 3^{100}$ par $11$. \hfill (0,5 pt)
	\item
	\begin{enumerate}
		\item Donner une solution particulière de l'équation : $23x - 17y = 1$. \hfill (0,25 pt)
		\item En déduire la résolution dans $\Z \times \Z$ de l'équation : $23x - 17y = 2$. \hfill (0,75 pt)
	\end{enumerate}
\end{enumerate}


% ------------------------------------------------------------
%   EXERCICE 2
% ------------------------------------------------------------

\exercice{2 points}

Dans l'espace muni du repère orthonormé direct $(O,\vec{\imath},\vec{\jmath},\vec{k})$.

\begin{enumerate}
	\item Écrire une équation paramétrique de la droite $(D)$ passant par le point $A(2\,;-1\,;3)$ et de vecteur directeur $\vec{u}(3\,;-2\,;1)$. \hfill (0,25 pt)
	\item Soit $(D')$ la droite d'équation paramétrique :
	$$
	\begin{cases}
		x = 4t+1 \\
		y = 5t-8 \\
		z = -2t+6
	\end{cases}
	\quad t \in \R.
	$$
	\begin{enumerate}
		\item Montrer que $(D)$ et $(D')$ sont orthogonaux. \hfill (0,75 pt)
		\item Déterminer les coordonnées du point $I$ intersection de $(D)$ et $(D')$. \hfill (0,25 pt)
	\end{enumerate}
	\item Écrire une équation cartésienne du plan $(P)$ contenant les deux droites $(D)$ et $(D')$. \hfill (0,75 pt)
\end{enumerate}


% ------------------------------------------------------------
%   PROBLEME 1
% ------------------------------------------------------------

\probleme{7 points}

Le plan $(P)$ étant orienté. On considère un carré direct $ABCD$ de centre $O$ tel que $AB = 4\text{cm}$.

$I$ étant le point défini par $\vect{AI} = \dfrac{3}{4}\vect{AB}$. Soit $J$ un point tel que $AJO$ est un triangle rectangle en $A$ avec $(\vect{AJ},\vect{AO}) = \dfrac{\pi}{2}$ et $AJ = 4\sqrt{2}\text{cm}$.

On note par :
\begin{itemize}[label=\textbullet]
	\item $r$ : la rotation de centre $O$ et d'angle $\dfrac{\pi}{2}$
	\item $h$ : l'homothétie de centre $I$ et de rapport $\sqrt{2}$
\end{itemize}

On pose $S = h \circ r$.

\medskip
\noindent\textbf{Partie A}

\begin{enumerate}
	\item Faire une figure et construire les points $I$ et $J$. \hfill (0,75 pt)
	\item
	\begin{enumerate}
		\item Déterminer et construire le point $G = bar\{(J,1)\,;(O,-4)\}$. \hfill (0,25 pt $\times$ 2)
		\item Déterminer et construire l'ensemble $(\Gamma) = \{M \in (P) \mid \left\Vert \vect{MJ} \right\Vert^2 - 4\left\Vert \vect{MO} \right\Vert^2 = 0\}$. \hfill (0,25 pt $\times$ 2)
	\end{enumerate}
	\item Déterminer et construire l'ensemble $(\Gamma') = \{M \in (P) \mid \vect{MO} \cdot \vect{MJ} = 0\}$. \hfill (0,25 pt $\times$ 2)
	\item
	\begin{enumerate}
		\item Déterminer le rapport et un angle de $S$. \hfill (0,25 pt $\times$ 2)
		\item Calculer $S(A)$. Conclure. \hfill (0,25 pt $\times$ 2)
		\item Montrer que le point $A$ est un point d'intersection de $(\Gamma)$ et $(\Gamma')$. \hfill (0,25 pt + 0,25 pt)
	\end{enumerate}
\end{enumerate}

\medskip
\noindent\textbf{Partie B}

Le plan $(P)$ est maintenant muni d'un repère orthonormé direct $(A,\vec{u},\vec{v})$ tel que $\vec{u} = \dfrac{1}{2}\vect{AB}$.

\begin{enumerate}
	\item Déterminer l'affixe des points $O$, $I$, $J$ et $G$. \hfill (0,25 pt $\times$ 4)
	\item
	\begin{enumerate}
		\item Déterminer l'expression complexe de $r$ et $h$. \hfill (0,25 pt $\times$ 4)
		\item En déduire l'expression complexe, la nature et les éléments caractéristiques de $S$. \hfill (0,25 pt $\times$ 2)
	\end{enumerate}
	\item Montrer que le point $A$ est un point d'intersection de $(\Gamma)$ et $(\Gamma')$. \hfill (0,75 pt)
\end{enumerate}


% ------------------------------------------------------------
%   PROBLEME 2
% ------------------------------------------------------------

\probleme{8 points}

Soit $f$ la fonction numérique définie sur $\R$ par :
$$
f(x) =
\begin{cases}
	(x+1)\e^{-x} - 1, & \text{si } x \leqslant 0 \\
	x(1-\ln x), & \text{si } x > 0
\end{cases}
$$
Notons par $(\mathscr{C})$ sa courbe représentative dans un repère orthonormé $(O,\vec{\imath},\vec{\jmath})$ d'unité $1\text{cm}$.

\medskip
\noindent\textbf{Partie I}

\begin{enumerate}
	\item Étudier la continuité et la dérivabilité de $f$ en $x_0 = 0$. \hfill (0,25 pt + 0,75 pt)
	\item
	\begin{enumerate}
		\item Étudier les variations de $f$ sur $\left]-\infty\,;0\right]$ et sur $\left]0\,;+\infty\right[$. \hfill (0,5 pt $\times$ 2)
		\item Dresser le tableau de variation de $f$ sur $\left]-\infty\,;+\infty\right[$. \hfill (1 pt)
	\end{enumerate}
	\item Étudier les branches infinies de $(\mathscr{C})$. \hfill (0,25 pt $\times$ 2)
	\item Construire la courbe $(\mathscr{C})$ en précisant la tangente ou les demi-tangentes au point d'abscisse $x_0 = 0$. \hfill (1,25 pt)
\end{enumerate}

\medskip
\noindent\textbf{Partie II}

Soit $(U_n)_{n\in\N}$ la suite définie par : $\forall n \in \N^*$ : $U_n = \dfrac{1}{n!}\displaystyle\int_{-1}^{0} (x+1)^n \e^{-x}\,dx$.

\begin{enumerate}
	\item
	\begin{enumerate}
		\item Calculer $U_1$. \hfill (0,25 pt)
		\item Interpréter géométriquement ce résultat. \hfill (0,25 pt)
	\end{enumerate}
	\item
	\begin{enumerate}
		\item En utilisant un encadrement de $\e^{-x}$ sur l'intervalle $[-1\,;0]$, donner un encadrement de $U_n$. \hfill (0,25 pt)
		\item En déduire $\displaystyle\lim_{n\to+\infty} U_n$. \hfill (0,25 pt)
	\end{enumerate}
	\item En utilisant une intégration par parties, vérifier que $\forall n \in \N^*$ : $U_{n+1} = U_n - \dfrac{1}{(n+1)!}$. \hfill (1,5 pt)
	\item Montrer que $\forall n \in \N^*$ : $\displaystyle\lim_{n\to+\infty} \sum_{k=0}^{n} \frac{1}{k!} = \e$. \hfill (0,75 pt)
\end{enumerate}
